% Algebra 2 Lesson 4-4 Lesson Quiz, short form (from the 2024 TE "4-4 Lesson Quiz" block, items 3-5).
% Items keep their Savvas numbers so the key, the packet's "On the quiz" box and lesson-prep.json agree.
\documentclass[11pt]{article}
\usepackage[margin=0.7in]{geometry}
\usepackage{amsmath,amssymb}
\usepackage{enumitem}
\usepackage{tikz}
\usepackage{pgfplots}
\pgfplotsset{compat=1.18}
\usepackage{fancyhdr}
\usepackage{parskip}
\setlist[enumerate]{leftmargin=*, itemsep=6pt}
\pagestyle{fancy}
\fancyhf{}
\fancyfoot[L]{\small\textit{en}Vision\textsuperscript{\textregistered} Algebra 2}
\fancyfoot[C]{\small 4-4 Lesson Quiz, short form}
\fancyfoot[R]{\small\thepage}
\renewcommand{\headrulewidth}{0pt}
\newcommand{\choice}[1]{\tikz[baseline=-0.6ex]\node[draw, circle, inner sep=0.6pt, minimum size=3.3mm, font=\scriptsize\bfseries] {#1};\;}
\newcommand{\checkbox}{$\square$\;}
\pgfplotsset{quizgrid/.style={
  axis lines=middle, grid=both, grid style={gray!45},
  ticklabel style={font=\scriptsize, fill=white, inner sep=1pt},
  xlabel={$x$}, ylabel={$y$}, axis line style={thick}, clip=true}}

\begin{document}

\noindent\textbf{Name}\; \rule{6cm}{0.4pt} \hfill Period \rule{1.2cm}{0.4pt} \quad Date \rule{2.2cm}{0.4pt}

\vspace{4pt}\noindent\rule{\textwidth}{0.6pt}

\vspace{2pt}\noindent{\Large\textbf{4-4 \; Lesson Quiz: Adding and Subtracting Rational Expressions}}

\noindent{\small short form: questions 3--5 of the Savvas 4-4 Lesson Quiz. Show your work for every question.}

\vspace{2pt}\noindent\rule{\textwidth}{0.6pt}

\vspace{6pt}
\begin{minipage}[t]{0.50\textwidth}
\textbf{3.} Subtract $\dfrac{3}{x-1}-\dfrac{4}{x+3}$.

\vspace{6pt}
\begin{itemize}[leftmargin=1.2em, itemsep=7pt]
\item[] \choice{A} $\dfrac{-x+13}{(x-1)(x+3)}$; $x\neq-3$ or $1$
\item[] \choice{B} $\dfrac{-x+5}{(x-1)(x+3)}$; $x\neq-3$ or $1$
\item[] \choice{C} $\dfrac{7x+5}{(x-1)(x+3)}$; $x\neq-3$ or $1$
\item[] \choice{D} $\dfrac{7x+13}{(x-1)(x+3)}$; $x\neq-3$ or $1$
\end{itemize}
\end{minipage}\hfill
\begin{minipage}[t]{0.47\textwidth}
\textbf{Work}

\vspace{4pt}
\foreach \i in {1,...,5}{\rule{0pt}{22pt}\rule{\linewidth}{0.4pt}\par}
\end{minipage}

\vspace{6pt}\noindent\rule{\textwidth}{0.4pt}

\vspace{6pt}
\begin{minipage}[t]{0.50\textwidth}
\textbf{4.} Simplify $\dfrac{\frac{1}{x}-\frac{1}{x+2}}{2}$.

\vspace{6pt}
\begin{itemize}[leftmargin=1.2em, itemsep=7pt]
\item[] \choice{A} $\dfrac{1}{x+2}$; $x\neq-2$
\item[] \choice{B} $\dfrac{1}{x}$; $x\neq0$
\item[] \choice{C} $\dfrac{1}{x(x+2)}$; $x\neq-2$ or $0$
\item[] \choice{D} $\dfrac{1}{x(x-2)}$; $x\neq2$ or $0$
\end{itemize}
\end{minipage}\hfill
\begin{minipage}[t]{0.47\textwidth}
\textbf{Work}

\vspace{4pt}
\foreach \i in {1,...,5}{\rule{0pt}{22pt}\rule{\linewidth}{0.4pt}\par}
\end{minipage}

\vspace{6pt}\noindent\rule{\textwidth}{0.4pt}

\vspace{6pt}
\textbf{5.} The equation $r=\dfrac{1}{\frac{1}{r_1}+\frac{1}{r_2}}$ represents the total resistance, $r$, when two resistors whose resistances are $r_1$ and $r_2$ are connected in parallel. Find the total resistance when $r_1$ is $x$ and $r_2$ is $x+1$.

\vspace{4pt}
\begin{minipage}[t]{0.50\textwidth}
\begin{itemize}[leftmargin=1.2em, itemsep=7pt]
\item[] \choice{A} $\dfrac{1}{2x+1}$; $x\neq-1,\,-\frac12$, or $0$
\item[] \choice{B} $2x+1$; $x\neq-1$ or $0$
\item[] \choice{C} $\dfrac{x(x+1)}{2x+1}$; $x\neq-1,\,-\frac12$, or $0$
\item[] \choice{D} $\dfrac{2x+1}{x(x+1)}$; $x\neq-1$ or $0$
\end{itemize}
\end{minipage}\hfill
\begin{minipage}[t]{0.47\textwidth}
\textbf{Work}

\vspace{4pt}
\foreach \i in {1,...,5}{\rule{0pt}{22pt}\rule{\linewidth}{0.4pt}\par}
\end{minipage}

\vfill
\noindent\textbf{Scoring:} Questions 3 and 4 are 10 points each; Question 5 is a bonus of up to 5. Total out of 20.

\newpage
\noindent{\LARGE\textbf{4-4 Lesson Quiz, short form: answer key}}

\vspace{4pt}\noindent\rule{\textwidth}{0.6pt}

\begin{enumerate}
\item[\textbf{3.}] \textbf{A, $\frac{-x+13}{(x-1)(x+3)}$; $x\neq-3$ or $1$.} LCD $(x-1)(x+3)$: $\frac{3(x+3)-4(x-1)}{(x-1)(x+3)}=\frac{3x+9-4x+4}{(x-1)(x+3)}=\frac{-x+13}{(x-1)(x+3)}$. Distractor B ($-x+5$) writes $-4(x-1)$ as $-4x-4$: the negative was not distributed to the $-1$. C ($7x+5$) adds the numerators instead of subtracting: $3(x+3)+4(x-1)=3x+9+4x-4$. D ($7x+13$) adds and also writes $4(x-1)$ as $4x+4$: $3x+9+4x+4$. Savvas quiz item 3, DOK 1, skill A86; taught by Example 4. \textbf{10 points:} 3 for the LCD, 4 for the numerator with the $-1$ distributed, 3 for choice A (domain included).
\item[\textbf{4.}] \textbf{C, $\frac{1}{x(x+2)}$; $x\neq-2$ or $0$.} Multiply top and bottom by the LCD $x(x+2)$: $\frac{(x+2)-x}{2x(x+2)}=\frac{2}{2x(x+2)}=\frac{1}{x(x+2)}$. Distractor D has the wrong sign in the factor $x-2$. Savvas quiz item 4, DOK 2, skill A86; taught by Example 6. \textbf{10 points:} 3 for the LCD $x(x+2)$, 4 for clearing the small fractions to $\frac{2}{2x(x+2)}$, 3 for choice C.
\item[\textbf{5.}] \textbf{C, $\frac{x(x+1)}{2x+1}$; $x\neq-1,\,-\frac12$, or $0$.} Substitute: $r=\frac{1}{\frac1x+\frac1{x+1}}$. Multiply top and bottom by $x(x+1)$: \mbox{$\frac{x(x+1)}{(x+1)+x}=\frac{x(x+1)}{2x+1}$}. $x=-\frac12$ is excluded because it makes $\frac1x+\frac1{x+1}=0$, so $r$ is undefined there. Distractor D is the reciprocal (the sum $\frac1x+\frac1{x+1}$ before dividing). Savvas quiz item 5, DOK 2, skill A86; Example 6 (the packet's Practice block is its parallel). \textbf{Bonus 5:} 2 for the substitution, 3 for choice C with work.
\end{enumerate}

\textbf{Items left off this form} (available for a retake or Form B): Savvas item 1, add $\frac{5}{2x+1}+\frac{8}{4x+2}$, $x\neq-\frac12$ (answer A, $\frac{9}{2x+1}$; DOK 1, A86; Example 3 content, practised in Try It 4a and IXL FEX); Savvas item 2, LCM of $(x+3)^2(x-2)$ and $(x+3)(x^2-16)$ (answer D, $(x+3)^2(x-2)(x^2-16)$; DOK 2, A86; Example 2 content, step 1 of Example 4).

\textbf{IXL replacement:} SmartScore 70 or above on FEX and YD8 replaces this quiz with 20/20.

\end{document}
